\subsection*{A complement bound without a positive spectral gap}
The ordinary operator residual permits a different reduction from the
coercive Schur estimate.  It does not require a positive gap on the
complement and therefore introduces no division by a collapsing level.

\begin{proposition}[A gap-free near-radical block bound]
\label{prop:v118-gap-free}
Let $A$ be a lower-bounded self-adjoint operator with closed form $q$,
and let $P$ be a finite-rank orthogonal projection with
$\operatorname{Ran}P\subset\mathcal D(A)$.  Put $Q=I-P$ and assume
$\norm{AP}\le\epsilon$.  Then $Q$ preserves the form domain and
\begin{equation}\label{eq:v118-gap-free-difference}
 \left|q(f,f)-q(Qf,Qf)\right|
 \le\frac{2}{\sqrt3}\epsilon\norm f^2
 \qquad(f\in\mathcal D(q)).
\end{equation}
Consequently, if
\[
 q(g,g)\ge-\eta\norm g^2\quad
 (g\in\operatorname{Ran}Q\cap\mathcal D(q)),\qquad \eta\ge0,
\]
then
\begin{equation}\label{eq:v118-gap-free-lower}
 \inf\sigma(A)\ge-\eta-\frac{2}{\sqrt3}\epsilon.
\end{equation}
The constant $2/\sqrt3$ in both statements is sharp.
\end{proposition}
\begin{proof}
Write $h=Pf$ and $g=Qf$.  The vector $h$ lies in the operator domain,
so $g$ belongs to the form domain and, with the fixed first-slot-linear
convention,
\[
 q(f,f)-q(g,g)
 =\ip{Ah}{h}+2\Re\ip{Ah}{g}
 =\Re\ip{Ah}{h+2g}.
\]
Since $h\perp g$, the absolute value is at most
\[
 \epsilon\norm h\sqrt{\norm h^2+4\norm g^2}
 \le\frac{2}{\sqrt3}\epsilon
           (\norm h^2+\norm g^2).
\]
For the last inequality, the maximum of
$t\sqrt{4-3t^2}$ on $0\le t\le1$ is $2/\sqrt3$.
This proves \eqref{eq:v118-gap-free-difference}; the complementary
lower bound and $\norm{Qf}\le\norm f$ prove
\eqref{eq:v118-gap-free-lower}.  Sharpness already holds on
$\mathbb C^2$: take $P$ onto the first coordinate and
\[
 A=-\frac{\epsilon}{\sqrt3}
   \begin{pmatrix}1&\sqrt2\\\sqrt2&0\end{pmatrix}.
\]
Then $\norm{AP}=\epsilon$, the complementary form is zero, and the
lowest eigenvalue is $-2\epsilon/\sqrt3$.
\end{proof}

For the rank-one projection onto either normalized source of
Proposition~\ref{prop:v115-absolute-residual},
$\norm{\mathsf W_\lambda P}=\norm{\mathsf W_\lambda u_\lambda}$
already tends to zero with the proved ordinary residual bound.
It therefore suffices to prove asymptotic nonnegativity
$q_\lambda|_{u_\lambda^\perp}\succeq-\eta_\lambda I$ with
$\eta_\lambda\to0$ along an unbounded window sequence.  Strict
complementary coercivity is unnecessary for this version of the
reduction.  For a growing source block the required uniform
$\norm{\mathsf W_\lambda P_\lambda}\to0$ remains a separate open
input.  Neither complementary sign estimate is proved here.

The missing sign cannot be concealed by removing the near-radical
block.  Indeed, suppose a fixed unit compact test $f$ has
$q(f,f)=-\delta<0$, and extend it by zero to every larger window.
Then \eqref{eq:v118-gap-free-difference} gives
\[
 q_\lambda((I-P_\lambda)f,(I-P_\lambda)f)
 \le-\delta+\frac{2}{\sqrt3}\epsilon_\lambda.
\]
Whenever $\epsilon_\lambda<\sqrt3\delta/2$, the projected vector
is nonzero and negative; its normalized Rayleigh value is bounded
above by the same negative quantity, since its norm is at most one.
Thus a negative direction persists in the complement of any uniformly
near-radical block, regardless of its rank.  By the existing cofinal
Weil-positivity criterion, the desired asymptotic complementary sign
for the proved rank-one source remains RH-strength.  The reduction
removes a gap denominator; it does not supply the arithmetic sign.
